% Options for packages loaded elsewhere
\PassOptionsToPackage{unicode}{hyperref}
\PassOptionsToPackage{hyphens}{url}
\documentclass[
  ignorenonframetext,
]{beamer}
\newif\ifbibliography
\usepackage{pgfpages}
\setbeamertemplate{caption}[numbered]
\setbeamertemplate{caption label separator}{: }
\setbeamercolor{caption name}{fg=normal text.fg}
\beamertemplatenavigationsymbolsempty
% remove section numbering
\setbeamertemplate{part page}{
  \centering
  \begin{beamercolorbox}[sep=16pt,center]{part title}
    \usebeamerfont{part title}\insertpart\par
  \end{beamercolorbox}
}
\setbeamertemplate{section page}{
  \centering
  \begin{beamercolorbox}[sep=12pt,center]{section title}
    \usebeamerfont{section title}\insertsection\par
  \end{beamercolorbox}
}
\setbeamertemplate{subsection page}{
  \centering
  \begin{beamercolorbox}[sep=8pt,center]{subsection title}
    \usebeamerfont{subsection title}\insertsubsection\par
  \end{beamercolorbox}
}
% Prevent slide breaks in the middle of a paragraph
\widowpenalties 1 10000
\raggedbottom
\AtBeginPart{
  \frame{\partpage}
}
\AtBeginSection{
  \ifbibliography
  \else
    \frame{\sectionpage}
  \fi
}
\AtBeginSubsection{
  \frame{\subsectionpage}
}
\usepackage{iftex}
\ifPDFTeX
  \usepackage[T1]{fontenc}
  \usepackage[utf8]{inputenc}
  \usepackage{textcomp} % provide euro and other symbols
\else % if luatex or xetex
  \usepackage{unicode-math} % this also loads fontspec
  \defaultfontfeatures{Scale=MatchLowercase}
  \defaultfontfeatures[\rmfamily]{Ligatures=TeX,Scale=1}
\fi
\usepackage{lmodern}
\ifPDFTeX\else
  % xetex/luatex font selection
\fi
% Use upquote if available, for straight quotes in verbatim environments
\IfFileExists{upquote.sty}{\usepackage{upquote}}{}
\IfFileExists{microtype.sty}{% use microtype if available
  \usepackage[]{microtype}
  \UseMicrotypeSet[protrusion]{basicmath} % disable protrusion for tt fonts
}{}
\makeatletter
\@ifundefined{KOMAClassName}{% if non-KOMA class
  \IfFileExists{parskip.sty}{%
    \usepackage{parskip}
  }{% else
    \setlength{\parindent}{0pt}
    \setlength{\parskip}{6pt plus 2pt minus 1pt}}
}{% if KOMA class
  \KOMAoptions{parskip=half}}
\makeatother
\usepackage{longtable,booktabs,array}
\newcounter{none} % for unnumbered tables
\usepackage{calc} % for calculating minipage widths
\usepackage{caption}
% Make caption package work with longtable
\makeatletter
\def\fnum@table{\tablename~\thetable}
\makeatother
\setlength{\emergencystretch}{3em} % prevent overfull lines
\providecommand{\tightlist}{%
  \setlength{\itemsep}{0pt}\setlength{\parskip}{0pt}}
% Auto-generated by infrastructure/rendering/_pdf_latex_helpers.py::extract_math_font_preamble.
% Minimal header passed to Pandoc via -H so Beamer slides pick up the same math font
% as the combined-PDF path. Adjust the manuscript's preamble.md to change the font globally.
\usepackage{unicode-math}
\setmathfont{latinmodern-math.otf}

% Slide layout defaults for warning-clean scientific decks.
\usepackage{etoolbox}
\IfFileExists{xurl.sty}{\usepackage{xurl}}{}
\IfFileExists{seqsplit.sty}{\usepackage{seqsplit}}{\newcommand{\seqsplit}[1]{#1}}
\protected\def\breakseq#1{\seqsplit{#1}}
\protected\def\breaktt#1{\begingroup\ttfamily\seqsplit{#1}\endgroup}
\setlength{\emergencystretch}{6em}
\tolerance=5000
\hbadness=10000
\hfuzz=1pt
\setlength{\tabcolsep}{2pt}
\AtBeginEnvironment{longtable}{\tiny\renewcommand{\arraystretch}{0.86}\setlength{\tabcolsep}{1pt}}
\AtBeginEnvironment{tabular}{\tiny\renewcommand{\arraystretch}{0.86}\setlength{\tabcolsep}{1pt}}

% Natbib and cross-reference fallbacks — slides don't load natbib
% or cleveref, but combined-PDF manuscript prose may emit these
% commands. The fallback renders citations as a bracketed key list
% and cross-references as detokenized labels so slides don't fail on
% undefined control sequences or raw underscores. \providecommand is
% a no-op if packages load later.
\providecommand{\citep}[1]{[#1]}
\providecommand{\citet}[1]{#1}
\providecommand{\citealp}[1]{#1}
\providecommand{\citeauthor}[1]{#1}
\providecommand{\citeyear}[1]{#1}
\providecommand{\cref}[1]{\texttt{\detokenize{#1}}}
\providecommand{\Cref}[1]{\texttt{\detokenize{#1}}}
\usepackage{bookmark}
\IfFileExists{xurl.sty}{\usepackage{xurl}}{} % add URL line breaks if available
\urlstyle{same}
% fallback for those not using the hyperref driver hyperxmp:
\makeatletter
\@ifundefined{xmpquote}{\newcommand{\xmpquote}[1]{#1}}{}
\makeatother
\hypersetup{
  hidelinks,
  pdfcreator={LaTeX via pandoc}}

\author{\texorpdfstring{}{}}
\date{}

\begin{document}

\section{Experimental Setup}\label{sec:experimental_setup}

\begin{frame}[allowframebreaks]{Experimental Setup}
This section details the controlled vocabulary, worked examples, and
software environment used to produce the results.
\end{frame}

\begin{frame}[fragile,allowframebreaks]{Controlled vocabulary}
\protect\phantomsection\label{controlled-vocabulary}
The DSL's vocabulary is declared once in code and consulted by every
gate and the compiler --- never re-declared per method:

{\def\LTcaptype{none} % do not increment counter
\begin{longtable}[]{@{}
  >{\raggedright\arraybackslash}p{(\linewidth - 4\tabcolsep) * \real{0.3333}}
  >{\raggedright\arraybackslash}p{(\linewidth - 4\tabcolsep) * \real{0.3333}}
  >{\raggedright\arraybackslash}p{(\linewidth - 4\tabcolsep) * \real{0.3333}}@{}}
\toprule\noalign{}
\begin{minipage}[b]{\linewidth}\raggedright
Module
\end{minipage} & \begin{minipage}[b]{\linewidth}\raggedright
Declares
\end{minipage} & \begin{minipage}[b]{\linewidth}\raggedright
Cardinality
\end{minipage} \\
\midrule\noalign{}
\endhead
\breaktt{src/methods\_dsl/vocabulary.py} & \texttt{StepKind},
\texttt{Target}, \texttt{target\_accepts} & 9 step kinds, 3 targets \\
\breaktt{src/methods\_dsl/units.py} & \texttt{Dimension},
\texttt{Quantity}, the unit table & 18 controlled units across 7
dimensions \\
\breaktt{src/methods\_dsl/validation.py} & The four staged gates & 4
gates, fixed order \\
\bottomrule\noalign{}
\end{longtable}
}
\end{frame}

\begin{frame}[fragile,allowframebreaks]{Worked examples}
\protect\phantomsection\label{worked-examples}
\breaktt{all\_example\_methods()}
(\breaktt{src/methods\_dsl/examples\_methods.py}) returns 2 methods:

{\def\LTcaptype{none} % do not increment counter
\begin{longtable}[]{@{}
  >{\raggedright\arraybackslash}p{(\linewidth - 6\tabcolsep) * \real{0.2500}}
  >{\raggedright\arraybackslash}p{(\linewidth - 6\tabcolsep) * \real{0.2500}}
  >{\raggedright\arraybackslash}p{(\linewidth - 6\tabcolsep) * \real{0.2500}}
  >{\raggedright\arraybackslash}p{(\linewidth - 6\tabcolsep) * \real{0.2500}}@{}}
\toprule\noalign{}
\begin{minipage}[b]{\linewidth}\raggedright
Method
\end{minipage} & \begin{minipage}[b]{\linewidth}\raggedright
Domain
\end{minipage} & \begin{minipage}[b]{\linewidth}\raggedright
Target
\end{minipage} & \begin{minipage}[b]{\linewidth}\raggedright
Notable structure
\end{minipage} \\
\midrule\noalign{}
\endhead
\texttt{PBSPreparation} & Wet-lab bench preparation (BPL's own domain;
an original example) & \texttt{HUMAN} & A strict 5-step linear chain
with a final \texttt{VALIDATE} step \\
\breaktt{SensorCalibrationSweep} & Instrument calibration (a non-biology
controlled procedure) & \texttt{AUTOMATED} & Mixed automated
\texttt{MEASURE}/\texttt{COMPUTE} steps and a \texttt{HUMAN}
\texttt{ANNOTATE} sign-off step \\
\bottomrule\noalign{}
\end{longtable}
}

The second example exists specifically to demonstrate that the
controlled vocabulary generalizes beyond wet-lab protocols, as
\breakseq{[@sec:introduction]} claims.
\end{frame}

\begin{frame}[fragile,allowframebreaks]{Pipeline conditions}
\protect\phantomsection\label{pipeline-conditions}
The experiment overlay (\breaktt{experiment\_plan.yaml}) declares three
conditions:

\begin{itemize}
\tightlist
\item
  \textbf{declared\_method} (reference) --- a method constructed
  directly as Python objects, before any validation gate has run.
\item
  \textbf{validated\_method} (proposed) --- the same method after
  passing all four staged gates and deterministic compilation to a
  scheduled \texttt{Plan}.
\item
  \textbf{automated\_target\_variant} (variant) ---
  \texttt{PBSPreparation}'s steps recompiled against an
  \texttt{AUTOMATED} execution target, ablating the target-compatibility
  gate's \texttt{HUMAN}/\texttt{AUTOMATED} step boundary (a
  \texttt{HUMAN} method's steps must all be \texttt{HUMAN}-compatible;
  an \texttt{AUTOMATED} method's steps may be either).
\end{itemize}

The primary metric is gate pass rate: the fraction of staged-gate
evaluations a method's steps satisfy.
\end{frame}

\begin{frame}[fragile,allowframebreaks]{Computational environment}
\protect\phantomsection\label{computational-environment}
\begin{itemize}
\tightlist
\item
  \textbf{Language}: Python 3.12.13 on Darwin arm64 (see root
  \texttt{pyproject.toml} for the supported version range).
\item
  \textbf{Core dependencies}: \texttt{pyyaml}, \texttt{matplotlib}
  (declared in \breaktt{domain\_profile.yaml::required\_packages}); the
  DSL library itself (\breaktt{src/methods\_dsl/}) has zero third-party
  dependencies beyond the standard library, with one declared
  \texttt{infrastructure} logging exception (\texttt{\_logging.py}).
\item
  \textbf{Headless plotting}: the analysis script sets
  \texttt{MPLBACKEND=Agg} before importing matplotlib.
\end{itemize}
\end{frame}

\begin{frame}[fragile,allowframebreaks]{Pipeline ordering}
\protect\phantomsection\label{pipeline-ordering}
The typical analysis order is:

\begin{enumerate}
\tightlist
\item
  \breaktt{scripts/methods\_analysis.py} --- compiles every example
  method, runs all gates, exports worklist/CSV/Mermaid/JSON per method,
  demonstrates the provenance hash-chain, and writes
  \breaktt{output/figures/step\_counts.png}, printing each output path
  for manifest collection.
\item
  \breaktt{scripts/z\_generate\_manuscript\_variables.py} --- reads
  \breaktt{manuscript/config.yaml} and the analysis outputs, then
  resolves every generated variable in \texttt{manuscript/*.md}.
\item
  PDF rendering reads the resolved manuscript tree so figure paths and
  prose match the analysis that just completed.
\end{enumerate}
\end{frame}

\begin{frame}[fragile,allowframebreaks]{Relation to results}
\protect\phantomsection\label{relation-to-results}
{\def\LTcaptype{none} % do not increment counter
\begin{longtable}[]{@{}
  >{\raggedright\arraybackslash}p{(\linewidth - 4\tabcolsep) * \real{0.3333}}
  >{\raggedright\arraybackslash}p{(\linewidth - 4\tabcolsep) * \real{0.3333}}
  >{\raggedright\arraybackslash}p{(\linewidth - 4\tabcolsep) * \real{0.3333}}@{}}
\toprule\noalign{}
\begin{minipage}[b]{\linewidth}\raggedright
Result (\breakseq{[@sec:results]})
\end{minipage} & \begin{minipage}[b]{\linewidth}\raggedright
Producing function (\breaktt{src/methods\_dsl/})
\end{minipage} & \begin{minipage}[b]{\linewidth}\raggedright
Primary inputs
\end{minipage} \\
\midrule\noalign{}
\endhead
Compiled-plan summary & \breaktt{compile\_method()} &
\breaktt{all\_example\_methods()} \\
Step-count figure & \texttt{len(plan.steps)} per method &
\breaktt{output/data/compiled\_plans.json} \\
Gate pass tally & \texttt{run\_all\_gates()} & Each example method \\
Determinism check & \breaktt{compile\_method()} called twice &
\breaktt{all\_example\_methods()} \\
Trust-chain verification & \texttt{append\_record()} /
\texttt{verify\_chain()} & Demonstration chain in
\breaktt{scripts/methods\_analysis.py} \\
\bottomrule\noalign{}
\end{longtable}
}

This table is descriptive documentation only; it is not executed as code
during the build.
\end{frame}

\end{document}
